% ============================================================
%  Objetivo 2 — InterPro y Pfam
% ============================================================

% ============================================================
\section{Ejercicio 2a}

\begin{consigna}
En la solapa Overview de InterPro: ¿qué dominios se encontraron? ¿En qué bases de datos?
¿Son los mismos que los reportados por NCBI-CDD?
\end{consigna}

InterPro encontró \textbf{5 entries} para la Proteína X:

\begin{table}[H]
\centering
\footnotesize
\begin{tabular}{llP{4.2cm}P{2.8cm}}
\toprule
Accession & Short name & Nombre completo & Región \\
\midrule
IPR004090 & Chemotax\_Me-accpt\_rcpt & Chemotaxis methyl-accepting receptor     & C-terminal (familia) \\
IPR004089 & MCPsignal\_dom           & MCP signalling domain                    & $\sim$270--617 aa \\
IPR038158 & H-NOX\_domain\_sf        & H-NOX domain superfamily                 & $\sim$1--190 aa \\
IPR024096 & NO\_sig/Golgi\_transp    & NO signalling/Golgi transport ligand-bd  & $\sim$1--190 aa \\
IPR011644 & Heme\_NO-bd              & Heme NO-binding                          & $\sim$1--190 aa \\
\bottomrule
\end{tabular}
\caption{Entries de InterPro encontradas para la Proteína X.}
\end{table}

Son los \textbf{mismos dos dominios} que reportó NCBI-CDD (HNOB + MCPsignal/Tar), pero
InterPro los desglosa en más entradas porque integra distintos niveles de clasificación: familia
específica, superfamilia estructural y superfamilia funcional. Las tres entries del N-terminal
(IPR011644, IPR038158 e IPR024096) describen el mismo dominio de unión a hemo desde ángulos
distintos.

Las bases de datos miembro que contribuyeron matches:

\begin{table}[H]
\centering
\footnotesize
\begin{tabular}{lP{2.5cm}P{5.5cm}}
\toprule
Base de datos & Matches & Dominio detectado \\
\midrule
\textbf{Pfam}             & 2 & HNOB (PF07730) + MCPsignal (PF00672) \\
\textbf{SUPERFAMILY}      & 2 & H-NOX superfamily + NO signalling superfamily \\
\textbf{CATH-Gene3D}      & 2 & Clasificaciones estructurales de ambos dominios \\
\textbf{SMART}            & 1 & MCPsignal \\
\textbf{PANTHER}          & 1 & Familia MCP general \\
\textbf{PROSITE profiles} & 1 & Perfil PS50111 \\
\bottomrule
\end{tabular}
\end{table}

La diferencia principal con NCBI-CDD es que InterPro agrega información funcional extra (términos
GO como \textit{heme binding} y \textit{chemotaxis}) e integra 7 bases de datos distintas, mientras
que CDD muestra directamente los 2 dominios. Pfam y SMART son las fuentes comunes a ambas.

% ============================================================
\section{Ejercicio 2b}

\begin{consigna}
Teniendo en cuenta la función del dominio HNOB, ¿qué aminoácido espera que esté muy conservado?
\end{consigna}

El dominio HNOB une un grupo hemo para detectar NO con afinidad femtomolar. En todas las
hemoproteínas conocidas, el hierro del hemo está coordinado axialmente por una \textbf{histidina
proximal} que actúa como ligando directo del hierro. Esta histidina es absolutamente indispensable:
si se muta, el hemo no se une y la proteína pierde completamente su función de sensor. La propia
descripción de Pfam para PF07730 lo dice explícitamente: \textit{``binds heme via a covalent linkage
to histidine''}. Esperamos entonces que la \textbf{histidina (H)} sea el residuo más conservado en
el alineamiento del dominio.

% ============================================================
\section{Ejercicio 2c}

\begin{consigna}
¿Encuentra el aminoácido muy conservado del punto 2b?
\end{consigna}

Sí. En el alineamiento seed de HNOB la histidina aparece prácticamente invariante en la
\textbf{posición 103} del modelo HMM. Su score de emisión en esa posición es \textbf{0.266},
que es el valor más bajo para His en todo el modelo --- los demás aminoácidos tienen scores de
entre 3.0 y 6.1 en la misma posición, es decir que tienen probabilidades muchísimo más bajas.
Esto confirma que la posición está prácticamente restringida a histidina, consistente con su
rol como ligando axial del hemo.

% ============================================================
\section{Ejercicio 2d}

\begin{consigna}
Elija del alineamiento 3 columnas muy conservadas, 2 poco conservadas y 2 posiciones con
gaps/inserciones, y justifique.
\end{consigna}

Tomamos los datos directamente del archivo del modelo HMM (PF07700).

\textbf{3 columnas muy conservadas} (score de emisión bajo $\Rightarrow$ alta probabilidad en esa posición):

\begin{table}[H]
\centering
\footnotesize
\begin{tabular}{P{2.5cm}P{2cm}P{2.5cm}P{6cm}}
\toprule
Posición HMM & AA & Score & Justificación \\
\midrule
132 & \textbf{Y} (Tyr) & 0.136 & Residuo aromático del núcleo hidrofóbico del dominio \\
116 & \textbf{P} (Pro) & 0.154 & Introduce un quiebre estructural en el loop que rodea el hemo \\
103 & \textbf{H} (His) & 0.266 & Histidina proximal; ligando axial del hierro del hemo \\
\bottomrule
\end{tabular}
\caption{Posiciones más conservadas en el alineamiento seed de HNOB.}
\end{table}

\textbf{2 columnas poco conservadas} (score mínimo alto $\Rightarrow$ ningún aminoácido domina):

\begin{table}[H]
\centering
\footnotesize
\begin{tabular}{P{2.5cm}P{2cm}P{2.5cm}P{6cm}}
\toprule
Posición HMM & Mejor AA & Score & Justificación \\
\midrule
71 & F & 2.375 & Posición de superficie; tolera muchos aminoácidos distintos \\
40 & T & 2.300 & Loop variable sin restricciones funcionales fuertes \\
\bottomrule
\end{tabular}
\end{table}

\textbf{2 posiciones con mayor probabilidad de inserción} (score de transición m$\to$i bajo):

\begin{table}[H]
\centering
\footnotesize
\begin{tabular}{P{2.5cm}P{2.5cm}P{7.5cm}}
\toprule
Posición HMM & Score m$\to$i & Justificación \\
\midrule
31  & 0.816 & Loop entre elementos de estructura secundaria; longitud variable entre organismos \\
103 & 0.887 & Zona alrededor del sitio de unión al hemo; los loops en esa región son variables \\
\bottomrule
\end{tabular}
\end{table}

% ============================================================
\section{Ejercicio 2e}

\begin{consigna}
¿Coincide la apreciación del alineamiento múltiple con el logo del HMM?
\end{consigna}

Sí, el logo del HMM es completamente consistente con lo que observamos en el alineamiento:

La \textbf{posición 103} aparece en el logo con una columna H (His) muy alta, alcanzando
$\sim$4.0 bits de contenido de información --- la más alta de esa región. La probabilidad de His
en esa posición es 0.766 (76.6\%), seguida por Y = 0.050 y F = 0.036, o sea que His es
ampliamente dominante. La ocupancia es 0.995, lo que significa que casi todas las secuencias
del seed tienen un residuo en esa posición (no un gap). Esto coincide exactamente con la His
casi invariante del alineamiento.

Las posiciones poco conservadas del alineamiento aparecen en el logo como columnas bajas con
mezcla de colores, lo que refleja que muchos aminoácidos son igualmente posibles. Y las filas
\textit{ins.prob} del logo muestran barras naranjas en las posiciones 31 y 103, confirmando
la mayor probabilidad de inserción en esas zonas.

% ============================================================
\section{Ejercicio 2f}

\begin{consigna}
¿Qué número de posición tiene la His del punto 2b?
\end{consigna}

La histidina proximal del hemo está en la \textbf{posición 103} del modelo HMM de HNOB (PF07700).

% ============================================================
\section{Ejercicio 2g}

\begin{consigna}
¿Qué valores de emisión tiene esa posición para los distintos aminoácidos? ¿Es lo que esperaba?
\end{consigna}

En HMMER los valores de emisión son log-odds negativos: un valor \textbf{bajo} significa alta
probabilidad (ese aminoácido es frecuente en esa posición), y un valor \textbf{alto} significa
baja probabilidad.

\begin{table}[H]
\centering
\footnotesize
\begin{tabular}{P{1.5cm}P{2cm}P{1.5cm}P{2cm}}
\toprule
AA & Score & AA & Score \\
\midrule
\textbf{H} & \textbf{0.266} $\leftarrow$ & P & 5.418 \\
Y & 3.004 & M & 5.498 \\
F & 3.324 & W & 4.955 \\
K & 4.223 & I & 5.010 \\
L & 4.245 & G & 4.901 \\
R & 4.386 & T & 4.875 \\
E & 4.503 & V & 4.838 \\
A & 4.582 & N & 4.699 \\
S & 4.577 & Q & 4.720 \\
D & 4.637 & C & 6.092 \\
\bottomrule
\end{tabular}
\caption{Valores de emisión del HMM para la posición 103 (His proximal del hemo).}
\end{table}

Es exactamente lo que esperábamos: His tiene el score más bajo de toda la fila (0.266) y todos
los demás aminoácidos tienen valores de 3.0 para arriba, con Cys en el extremo opuesto (6.092).
La diferencia es muy grande, lo que confirma que esta posición está prácticamente restringida a
histidina. Tiene sentido biológicamente: sin esa His el hemo no se une y la proteína pierde su
función completamente, así que cualquier mutación en esa posición sería letal para la función del
sensor.

% ============================================================
\section{Ejercicio 2h}

\begin{consigna}
¿Con qué otros dominios se puede combinar HNOB según las arquitecturas de Pfam?
\end{consigna}

Pfam reporta 134 arquitecturas distintas para HNOB. Las más frecuentes son:

\begin{table}[H]
\centering
\footnotesize
\begin{tabular}{P{1.5cm}P{5cm}P{5cm}}
\toprule
Proteínas & Arquitectura & Función \\
\midrule
5189 & HNOB --- HNOBA --- Guanilato/adenilato ciclasa & Guanilato ciclasa soluble (sGC) de animales \\
2848 & HNOB (solo)                                     & Sensor de NO bacteriano sin efector acoplado \\
374  & HNOB --- MCPsignal                              & Receptor de quimiotaxis (\textbf{nuestra proteína}) \\
357  & HNOB --- HNOBA                                  & Con dominio asociado \\
\bottomrule
\end{tabular}
\caption{Arquitecturas de dominio más frecuentes en proteínas con dominio HNOB.}
\end{table}

Lo interesante es que la Proteína X de \textit{C. tetani} pertenece a una arquitectura minoritaria
(374 proteínas) comparada con la sGC de animales que tiene 5189. La diferencia en el efector es
clave: en bacterias el dominio HNOB se combina con MCPsignal para hacer quimiotaxis, mientras
que en animales se combina con HNOBA y una ciclasa para producir cGMP como segundo mensajero.
Mismo dominio sensor, distinto efector según el reino.

% ============================================================
\section{Ejercicio 2i}

\begin{consigna}
¿Cuántas especies de \textit{Clostridium} tienen el dominio HNOB según Pfam? ¿Cuáles son las
más representadas?
\end{consigna}

En la sección de taxonomía de Pfam para HNOB (PF07700), encontramos \textbf{118 taxa} de
\textit{Clostridium} con al menos una proteína con este dominio. Las más representadas:

\begin{table}[H]
\centering
\footnotesize
\begin{tabular}{lP{3cm}}
\toprule
Especie & Proteínas con HNOB \\
\midrule
\textit{Clostridium botulinum}      & 22 \\
\textit{Clostridium beijerinckii}   & 10 \\
\textit{Clostridium sporogenes}     &  5 \\
\textit{Clostridium estertheticum}  &  4 \\
\textit{Clostridium tetani}         &  \textbf{4} \\
\textit{Clostridium intestinale}    &  3 \\
\textit{Clostridium kluyveri}       &  3 \\
\textit{Clostridium cochlearium}    &  3 \\
\textit{Clostridium perfringens}    &  1 \\
\bottomrule
\end{tabular}
\caption{Especies de \textit{Clostridium} con dominio HNOB según Pfam (extracto).}
\end{table}

\textit{C. tetani} tiene 4 proteínas con dominio HNOB, incluyendo la Proteína X que estamos
estudiando. \textit{C. botulinum} es el más representado con 22 proteínas, posiblemente porque
tiene más genomas secuenciados. El dominio está presente en prácticamente todo el género, lo
que sugiere que es una función ancestral en \textit{Clostridium}.
